\begin{abstract}
We investigate the ground exciton and biexciton states in a hard-wall,
strongly oblate semi-ellipsoidal GaAs quantum dot. Their energies are evaluated
variationally using correlated trial states constructed from analytical
single-particle wave functions, together with a ground-orbital-matched
importance transformation and adaptive quasi-Monte Carlo integration. A
noise-aware coordinate pattern search with two-budget verification is used to optimize
the variational parameters along intersecting axial and lateral geometry
scans. Because the effective-mass hard-wall model is intended to isolate
geometrical effects, the principal results are qualitative trends rather than
absolute predictions for a particular fabricated dot. The axial dimension
controls the absolute confinement energy much more strongly than the lateral
dimension, whereas enlargement along either direction weakens the biexciton
binding. The central binding estimate remains positive throughout the sampled
range but is not resolved from zero at the widest lateral geometry. Increasing
either semiaxis enhances the exciton coincidence overlap and oscillator
strength and therefore shortens the calculated radiative lifetimes. The
lateral dependence remains pronounced where the axial dependence approaches
saturation. These results distinguish the dominant axial control of the
confinement-energy scale from the important lateral control of weak
biexciton binding and radiative strength.
\end{abstract}

\noindent\textbf{Keywords:} semi-ellipsoidal quantum dot; exciton;
biexciton; binding energy; variational method; quasi-Monte Carlo integration;
radiative lifetime
